\subsection{Interpretation}
  \label{subsec:interpretation}

  This toy model illustrates how two observables may appear correlated because they
  represent different projections of a single underlying structure.
  In this sense, Bell-type correlations may be understood as consistency constraints
  imposed by a non-factorizable effective description, rather than as signatures of
  nonlocal causation.

  While highly simplified, the model captures the essential logical mechanism discussed
  in the main text.
  More elaborate realizations may be constructed by enriching the structure of
  $\Omega$ or the projection $\Pi$, but the origin of the correlations remains the same:
  the non-injective mapping between underlying configurations and observable outcomes.
